\clearpage
\fancyhead[L]{\small PRECESSION.}
\SourceLeaf{566}
\begin{center}\small Göttingische gelehrte Anzeigen. No. 19. Pages 188..189. 31 January 1814.\\[3pt]\rule{.72\linewidth}{.3pt}\end{center}

The Royal Academy of Sciences in Berlin had offered a prize for the most accurate determination of the magnitude of precession. It awarded the prize to Mr. \textsc{Bessel} for his work on the question. Indeed, no one could have undertaken this determination with greater success than Mr. \textsc{Bessel}, who has been working for six years on a broad analysis of \textsc{Bradley}'s observations. For this purpose he drew on 4585 determinations for stars, and we believe we are doing astronomers an important service by giving here his final result for a quantity they need every day.

\begin{center}\fontsize{8}{11}\selectfont\renewcommand{\arraystretch}{1.18}\begin{tabular}{@{}l@{\hspace{4pt}}l@{}}
Lunisolar precession $.\;.\;.\;.\;.\;.\;.\;.\;.\;.\;.\;.$ & $= 50''35330 - 0''0002435890\,(t-1800)$ \\
Observed general precession $.\;.\;.\;.\;.\;.$ & $50.18924 + \mathrm{c}.0002442966\,(t-1800)$ \\
Constant for precession in right ascension & $46.01058 + 0.0003590677\,(t-1800)$ \\
Constant for precession in declination $.\;.$ & $20.04966 - 0.0002135621\,(t-1800)$ \\
\end{tabular}\end{center}

Mr. \textsc{Bessel} found no clear trace of the phenomenon noted by \textsc{Triesnecker}, in which the motions at some points in the sky were said to indicate unequal precession.
